Huntington-Hill is U.S.\ apportionment law, and it is a good law.
Its geometric mean rounding rule implements a mathematically principled
notion of equal per-capita representation across all pairs of states,
and its divisor-method structure guarantees immunity to the
Alabama paradox, the population paradox, and the new-states paradox.
These properties were the primary technical rationale for Congress's
1941 adoption of the method and the Supreme Court's unanimous
upholding in \textit{Montana} (1992).

The \textsc{bisect-apportion} implementation provides practitioners
with an exact, SHA-verified, auditable computation of the
Huntington-Hill apportionment for 2020, 2010, and 2000 Census data.
Its integer arithmetic ensures that priority comparisons are exact
even in near-tie cases, and its CLI provides a simple interface
for redistricting pipelines to obtain the statutory seat counts
before drawing district boundaries.

Three caveats deserve emphasis.
First, HH's optimality under the population pair test (Theorem~\ref{thm:hh-optimal})
is specific to the \emph{relative} deviation criterion;
Webster is optimal under the \emph{absolute} deviation criterion.
The choice between HH and Webster is not a mathematical question
but a political one about which notion of fairness Congress preferred.
Second, the quota rule --- which Hamilton satisfies but HH does not ---
remains a point of academic debate.
The Balinski-Young impossibility theorem establishes that no method can
satisfy both the quota rule and paradox immunity, but does not tell us
which property Congress should prioritise.
Third, the Montana case confirms that courts will not second-guess
Congress's choice of apportionment method; the role of mathematical
analysis is to characterise the method's properties, not to
litigate its superiority.

Practitioners using \textsc{bisect} for redistricting can rely on
\textsc{bisect-apportion} as a correct and legally defensible
implementation of 2~U.S.C.\ \S2a.
The SHA verification provides the audit trail that commission and
court contexts require.
